\documentclass[11pt,a4paper]{article}

\usepackage[utf8]{inputenc}
\usepackage[T1]{fontenc}
\usepackage{lmodern}
\usepackage[ngerman]{babel}
\usepackage[margin=2.5cm]{geometry}
\usepackage{microtype}
\usepackage{amsmath,amssymb,amsthm}
\usepackage{booktabs}
\usepackage{array}
\usepackage{graphicx}
\usepackage{xcolor}
\usepackage{tikz}
\usetikzlibrary{arrows.meta,calc,positioning}
\usepackage[ruled,vlined,linesnumbered]{algorithm2e}
\usepackage{float}
\usepackage{hyperref}
\hypersetup{
  colorlinks=true,
  linkcolor=blue!50!black,
  citecolor=blue!50!black,
  urlcolor=blue!60!black,
  pdftitle={Compass16-LBM: Lattice-Boltzmann-inspirierte Schwarmkoordination},
  pdfauthor={Dogan Balban},
  pdfsubject={Schwarmkoordination, Lattice-Boltzmann, Embedded Systems},
  pdfkeywords={LBM, swarm coordination, integer arithmetic, embedded systems, compass discretization}
}
\usepackage{enumitem}

% Improve breaking of long German compound words & hyphenated terms
\hyphenation{Lattice-Boltz-mann-Me-tho-de Schwarm-ko-or-di-na-tion
  Kom-pass-dis-kre-ti-sie-rung Mi-kro-con-trol-lern Mi-kro-con-trol-ler
  Re-yn-olds-Floc-king Ge-schwin-dig-keits-mo-dell
  Ge-schwin-dig-keits-vek-tor Ge-schwin-dig-keits-klas-sen
  Be-we-gungs-ten-denz Be-we-gungs-im-puls Be-we-gungs-rich-tung
  Sicher-heits-ab-stand Wahr-neh-mungs-schicht Sta-bi-li-t\"ats-be-din-gung
  Ant-ip-aral-le-li-t\"at Hin-der-nis-be-hand-lung Po-si-tions-up-date
  Re-la-xa-tions-schritt Gleich-ge-wichts-ver-tei-lung}

\newtheorem{definition}{Definition}
\newtheorem{theorem}{Theorem}
\newtheorem{lemma}{Lemma}
\newtheorem{proposition}{Proposition}
\newtheorem{remark}{Bemerkung}

\newcommand{\feq}{f^{\mathrm{eq}}}
\newcommand{\compass}{\mathrm{C16}}
\newcommand{\floor}[1]{\lfloor #1 \rfloor}

\title{\textbf{Compass16-LBM}\\[6pt]
\large Lattice-Boltzmann-inspirierte Schwarmkoordination\\
mit ganzzahliger Kompassdiskretisierung\\[12pt]
\normalsize Mathematische Formulierung --- Arbeitsdokument v2}
\author{Dogan Balban\\[4pt]
\small\textit{Independent Research --- M\"arz 2026}}
\date{}

\begin{document}
\maketitle

\begin{abstract}
\noindent Wir formulieren ein diskretes Geschwindigkeitsmodell f\"ur die dezentrale Schwarmkoordination
auf Integer-Hardware, inspiriert von der Lattice-Boltzmann-Methode (LBM).
Statt der \"ublichen D2Q9-Diskretisierung (8~Richtungen + Ruhe) verwenden wir eine
\textit{Compass16}-Diskretisierung (16~Richtungen + Ruhe = 17~Werte pro Agent),
die direkt aus dem DSM-Framework (Balban, 2026) abgeleitet ist.
Die gesamte Formulierung --- Gleichgewichtsverteilung, Relaxation, Positionsupdate ---
verwendet ausschlie\ss lich ganzzahlige Arithmetik mit Festkomma-Skalierung $S=256$.
Zielanwendung: Kollisionsfreie Schwarmkoordination auf 8/16-Bit-Mikrocontrollern ohne FPU.
\end{abstract}

% ============================================================
\section{Motivation: Von LBM zu Schwarmkoordination}
% ============================================================

Die klassische Lattice-Boltzmann-Methode simuliert Fl\"ussigkeiten durch zwei lokale Operationen:
\begin{enumerate}[itemsep=2pt]
\item \textbf{Relaxation:} Verteilungswerte an jedem Gitterknoten n\"ahern sich dem lokalen Gleichgewicht an.
\item \textbf{Propagation:} Jeder Verteilungswert bewegt sich einen Schritt in seine diskrete Richtung.
\end{enumerate}
Aus diesen rein lokalen Regeln entsteht \textit{emergentes} globales Verhalten:
Die Navier-Stokes-Gleichungen werden makroskopisch exakt reproduziert,
obwohl kein Agent die Gleichungen kennt oder l\"ost.

\textbf{Kernidee:} Wir \"ubertragen dieses Prinzip auf mobile Agenten (Roboter, Drohnen).
Statt Fl\"ussigkeitsteilchen koordinieren wir Agenten, die sich flie\ss end und kollisionsfrei
bewegen --- ohne zentrale Steuerung, rein durch lokale Interaktion.

\medskip
\noindent\textbf{Begriffsklärung.}
Die LBM-Literatur verwendet Begriffe aus der Fluiddynamik (Dichte, Druck, Schallgeschwindigkeit).
In unserem Schwarmmodell haben diese Gr\"o\ss en eine andere Bedeutung.
Um Verwechslungen zu vermeiden, verwenden wir durchgehend schwarmspezifische Terminologie:

\begin{center}
\begin{tabular}{@{}lll@{}}
\toprule
\textbf{LBM-Original} & \textbf{Schwarm-Bedeutung} & \textbf{Symbol} \\
\midrule
Gitterknoten & Agentenposition & $\mathbf{x}_A$ \\
Population $f_i$ & Bewegungstendenz in Richtung $i$ & $f_k$ \\
Lokale Dichte $\rho$ & Aktivit\"atslevel des Agenten & $\rho$ \\
Makro-Geschwindigkeit $\mathbf{u}$ & Mittlere Bewegungsrichtung & $\mathbf{u}$ \\
Schallgeschwindigkeit $c_s^2$ & \textit{entf\"allt} --- ersetzt durch Reaktivit\"at & $\kappa$ \\
Druck & Interaktionsst\"arke (Nachbarabsto\ss ung) & --- \\
Viskosit\"at & Tr\"agheit der Richtungs\"anderung & $\propto (S/\omega - 1/2)$ \\
Bounce-Back an W\"anden & Ausweichen an Hindernissen & Richtungstausch \\
\bottomrule
\end{tabular}
\end{center}

% ============================================================
\section{Compass16-Geschwindigkeitssatz}
% ============================================================

\begin{definition}[Compass16-Geschwindigkeitsvektoren]
Der diskrete Geschwindigkeitssatz $\mathcal{C}_{17} = \{\mathbf{c}_0, \mathbf{c}_1, \ldots, \mathbf{c}_{16}\}$
besteht aus dem Ruhevektor und 16 Richtungsvektoren im Abstand von $22{,}5^\circ$:
\begin{equation}
\mathbf{c}_k = 
\begin{cases}
(0, 0) & k = 0 \quad \text{(Ruhe)} \\[4pt]
\displaystyle\Big(\floor{256\sin\big(\tfrac{(k-1)\pi}{8}\big)},\;\;
\floor{256\cos\big(\tfrac{(k-1)\pi}{8}\big)}\Big) & k = 1,\ldots,16
\end{cases}
\end{equation}
Die Skalierung $S = 256 = 2^8$ erm\"oglicht effiziente Shift-Operationen auf 8-Bit-Hardware.
Die Werte werden einmalig vorberechnet und als Konstanten im ROM gespeichert.
\end{definition}

\noindent Die konkreten Vektoren:

\begin{center}
\small
\begin{tabular}{@{}clrrcclrr@{}}
\toprule
$k$ & Richtung & $c_{k,x}$ & $c_{k,y}$ & \quad &
$k$ & Richtung & $c_{k,x}$ & $c_{k,y}$ \\
\midrule
0  & Ruhe & 0 & 0       &  & 9  & S   & 0 & $-256$ \\
1  & N    & 0 & 256     &  & 10 & SSW & $-98$  & $-237$ \\
2  & NNE  & 98 & 237    &  & 11 & SW  & $-181$ & $-181$ \\
3  & NE   & 181 & 181   &  & 12 & WSW & $-237$ & $-98$ \\
4  & ENE  & 237 & 98    &  & 13 & W   & $-256$ & 0 \\
5  & E    & 256 & 0     &  & 14 & WNW & $-237$ & 98 \\
6  & ESE  & 237 & $-98$ &  & 15 & NW  & $-181$ & 181 \\
7  & SE   & 181 & $-181$&  & 16 & NNW & $-98$  & 237 \\
8  & SSE  & 98 & $-237$ &  &    &     &        &     \\
\bottomrule
\end{tabular}
\end{center}

\begin{remark}
Die Werte $98 \approx 256\sin(22{,}5^\circ)$ und $237 \approx 256\cos(22{,}5^\circ)$
sind dieselben Konstanten, die in DSM als Schwellwerte $\tau_1 = 0{,}199 \approx 51/256$
und $\tau_2 = 0{,}668 \approx 171/256$ auftreten.
Die Verbindung: $\sin(22{,}5^\circ)/\cos(22{,}5^\circ) = \tan(22{,}5^\circ) \approx \tau_1/\tau_2$ (bis auf Rundung).
Der gesamte Compass16-Klassifikator aus DSM bleibt unver\"andert verwendbar.
\end{remark}

\begin{remark}[Antiparallelit\"at]
F\"ur jeden Richtungsindex $k \in \{1,\ldots,16\}$ ist der antiparallele Index:
\begin{equation}
\bar{k} = ((k - 1 + 8) \bmod 16) + 1
\end{equation}
Es gilt: $\mathbf{c}_k + \mathbf{c}_{\bar{k}} = \mathbf{0}$. Dies ist die DSM-Antiparallelit\"at,
die hier f\"ur die Hindernisbehandlung (Abschnitt~\ref{sec:bounce}) zentral wird.
\end{remark}

% ============================================================
\section{Zustandsvektor und makroskopische Gr\"o\ss en}
% ============================================================

\begin{definition}[Agentenzustand]
Jeder Agent $A$ speichert einen Verteilungsvektor:
\begin{equation}
\mathbf{f}^{(A)} = (f_0, f_1, f_2, \ldots, f_{16}) \in \mathbb{Z}_{\geq 0}^{17}
\end{equation}
Jeder Wert $f_k$ ist eine nicht-negative ganze Zahl und repr\"asentiert die
\textit{Bewegungstendenz} des Agenten in Richtung $\mathbf{c}_k$.
Ein hoher Wert $f_k$ bedeutet: Der Agent ``will'' sich stark in Richtung $k$ bewegen.
\end{definition}

\noindent Aus dem Verteilungsvektor leiten sich zwei makroskopische Gr\"o\ss en ab:

\textbf{Aktivit\"atslevel} (Gesamtintensit\"at aller Bewegungstendenzen):
\begin{equation}
\rho = \sum_{k=0}^{16} f_k
\end{equation}
Im station\"aren Zustand gilt $\rho \approx \rho_0$ (Referenzwert, typisch $\rho_0 = 256$).

\textbf{Bewegungsimpuls} (resultierende Richtung und Geschwindigkeit):
\begin{equation}
\mathbf{j} = \sum_{k=1}^{16} f_k \cdot \mathbf{c}_k
\end{equation}
Die mittlere Bewegungsrichtung ergibt sich als $\mathbf{u} = \mathbf{j} / \rho$.

In Integer-Arithmetik:
\begin{align}
j_x &= \sum_{k=1}^{16} f_k \cdot c_{k,x} \label{eq:jx}\\
j_y &= \sum_{k=1}^{16} f_k \cdot c_{k,y} \label{eq:jy}
\end{align}
Die Division $\mathbf{u} = \mathbf{j}/\rho$ wird durch Shift approximiert,
wenn $\rho$ im station\"aren Zustand auf $\rho_0 = 256$ normiert bleibt.

\textbf{Speicherbedarf pro Agent:} 17 Verteilungswerte $\times$ 16\,Bit = 34 Bytes.
Dazu: 2 Positionskoordinaten $\times$ 16\,Bit = 4 Bytes.
\textbf{Gesamt: 38 Bytes pro Agent.}

% ============================================================
\section{Gleichgewichtsverteilung}
% ============================================================

Die Gleichgewichtsverteilung $\feq_k$ beschreibt den Zustand, den ein Agent
\textit{anstreben soll}. Die Relaxation (Abschnitt~\ref{sec:collision}) treibt
den aktuellen Zustand $f_k$ kontinuierlich in Richtung $\feq_k$.

\begin{definition}[Compass16-Gleichgewicht]\label{def:feq}
\begin{equation}\label{eq:feq}
\feq_k = w_k \cdot \Big(\rho + \kappa \cdot (\mathbf{c}_k \cdot \mathbf{u}_{\mathrm{target}})\Big)
\end{equation}
wobei:
\begin{itemize}[itemsep=2pt]
\item $w_k$ = Richtungsgewicht (bestimmt die Isotropie, siehe Abschnitt~\ref{sec:weights})
\item $\rho$ = aktuelles Aktivit\"atslevel des Agenten
\item $\kappa \in \mathbb{Z}_{>0}$ = \textbf{Reaktivit\"at} (ganzzahliger Steuerparameter)
\item $\mathbf{u}_{\mathrm{target}}$ = \textbf{Sollbewegung} (zusammengesetzt aus Navigation, Separation, Koh\"asion)
\end{itemize}
\end{definition}

\begin{remark}[Bedeutung der Reaktivit\"at $\kappa$]
$\kappa$ steuert, wie stark die Sollbewegung $\mathbf{u}_{\mathrm{target}}$ die Verteilung beeinflusst:
\begin{itemize}[itemsep=2pt]
\item $\kappa = 1$: Minimale Reaktion. Agent \"andert seine Richtung nur z\"ogerlich.
Geeignet f\"ur tr\"age, schwere Roboter.
\item $\kappa = 4$: Moderate Reaktion. Guter Standardwert f\"ur Lagerroboter.
\item $\kappa = 16$: Starke Reaktion. Agent folgt der Sollbewegung fast instantan.
Geeignet f\"ur leichte, agile Einheiten.
\end{itemize}
Da $\kappa$ als Zweierpotenz gew\"ahlt werden kann ($1, 2, 4, 8, 16$),
ist die Multiplikation $\kappa \cdot (\mathbf{c}_k \cdot \mathbf{u}_{\mathrm{target}})$
als Links-Shift implementierbar: $\kappa = 2^n \;\Rightarrow\; \kappa \cdot x = x \ll n$.
\end{remark}

\subsection{Richtungsgewichte $w_k$}\label{sec:weights}

Die Gewichte m\"ussen die Symmetriebedingungen
$\sum w_k = 1$ und $\sum w_k c_{k,\alpha} c_{k,\beta} \propto \delta_{\alpha\beta}$
erf\"ullen, damit keine Richtung bevorzugt wird.

F\"ur 16 gleichm\"a\ss ig verteilte Richtungen plus Ruhe w\"ahlen wir:
\begin{equation}
w_k = \begin{cases}
w_0 = 1 - 16\,w_r & k = 0 \quad \text{(Ruhe)} \\
w_r = 1/24 & k = 1,\ldots,16 \quad \text{(alle Richtungen gleich)}
\end{cases}
\end{equation}

Dies ergibt $w_0 = 1/3$ und $w_r = 1/24$.

\textbf{Integer-Darstellung} mit Nenner-Skalierung $W = 768 = 3 \times 256$:
\begin{equation}
\hat{w}_0 = W \cdot w_0 = 256, \qquad \hat{w}_r = W \cdot w_r = 32
\end{equation}
Die Gleichgewichtsformel (\ref{eq:feq}) wird dann berechnet als:
\begin{equation}
\feq_k = \frac{\hat{w}_k \cdot \big(\rho + \kappa \cdot (\mathbf{c}_k \cdot \mathbf{u}_{\mathrm{target}}) \gg 8\big)}{W}
\end{equation}
wobei die Division durch $W = 768$ als $(\cdot \times 341) \gg 18$ approximiert werden kann
(da $341/2^{18} \approx 1/768$). Alternativ: Lookup-Table f\"ur die 17 Produkte $\hat{w}_k / W$.

\subsection{Sollbewegung $\mathbf{u}_{\mathrm{target}}$}

Die Sollbewegung setzt sich aus drei Komponenten zusammen,
die den Reynolds-Flocking-Regeln (1987) entsprechen:

\begin{equation}\label{eq:utarget}
\mathbf{u}_{\mathrm{target}} = \alpha\,\mathbf{u}_{\mathrm{Ziel}}
+ \beta\,\mathbf{u}_{\mathrm{Sep}}
+ \gamma\,\mathbf{u}_{\mathrm{Koh}}
\end{equation}

\textbf{(a) Navigation} (Attraktion zum Ziel):
\begin{equation}
\mathbf{u}_{\mathrm{Ziel}} = \mathbf{c}_{\kappa_Z}
\quad\text{wobei}\quad
\kappa_Z = \compass(\mathbf{x}_{\mathrm{Ziel}} - \mathbf{x}_A)
\end{equation}
Die Compass16-Klassifikation $\compass(\cdot)$ aus DSM: \textbf{keine Trigonometrie},
nur Schwellwertvergleiche mit $\tau_1 = 51/256$ und $\tau_2 = 171/256$.

\textbf{(b) Separation} (Absto\ss ung von zu nahen Nachbarn --- ``Gummiseil unter Druck''):
\begin{equation}
\mathbf{u}_{\mathrm{Sep}} = -\sum_{\substack{j \in \mathcal{N}(A) \\ d^2_j < d^2_{\min}}}
\frac{d^2_{\min}}{d^2_j} \cdot \mathbf{c}_{\compass(\mathbf{x}_j - \mathbf{x}_A)}
\end{equation}
Je n\"aher der Nachbar ($d^2_j \to 0$), desto st\"arker die Absto\ss ung.
Die Division $d^2_{\min}/d^2_j$ kann \"uber eine Lookup-Table (256 Eintr\"age, 512 Bytes ROM)
oder st\"uckweise lineare Approximation realisiert werden.

\textbf{(c) Koh\"asion} (Anziehung zu entfernten Nachbarn --- ``Gummiseil unter Zug''):
\begin{equation}
\mathbf{u}_{\mathrm{Koh}} = +\sum_{\substack{j \in \mathcal{N}(A) \\ d^2_j > d^2_{\max}}}
\frac{d^2_j}{d^2_{\max}} \cdot \mathbf{c}_{\compass(\mathbf{x}_j - \mathbf{x}_A)}
\end{equation}
Je weiter der Nachbar ($d^2_j \gg d^2_{\max}$), desto st\"arker der Zug.

Die Gewichte $\alpha, \beta, \gamma \in \mathbb{Z}$ sind ganzzahlige Steuerparameter.
Sie werden als Zweierpotenzen gew\"ahlt, um Multiplikationen durch Shifts zu ersetzen:

\begin{center}
\begin{tabular}{@{}llll@{}}
\toprule
\textbf{Parameter} & \textbf{Wert} & \textbf{Shift} & \textbf{Wirkung} \\
\midrule
$\alpha$ (Navigation) & 128 & $\ll 7$ & Zielgerichtete Bewegung \\
$\beta$ (Separation) & 256 & $\ll 8$ & Kollisionsvermeidung (st\"arkster Term) \\
$\gamma$ (Koh\"asion) & 64 & $\ll 6$ & Formationserhalt (schw\"achster Term) \\
\bottomrule
\end{tabular}
\end{center}

Dass $\beta > \alpha > \gamma$ gew\"ahlt wird, spiegelt die Priorit\"at wider:
Kollisionsvermeidung ist wichtiger als Zielnavigation, und Zielnavigation
ist wichtiger als Formationserhalt. Dies ist eine bewusste Entwurfsentscheidung.

% ============================================================
\section{Relaxationsschritt}\label{sec:collision}
% ============================================================

Im Relaxationsschritt n\"ahert sich der aktuelle Verteilungsvektor
dem Gleichgewicht an. Die Geschwindigkeit dieser Ann\"aherung
bestimmt, wie ``tr\"age'' oder ``reaktiv'' der Agent sich bewegt.

\begin{definition}[Integer-Relaxation (BGK-Schema)]
\begin{equation}\label{eq:collision}
f_k^* = f_k - \frac{\omega \cdot (f_k - \feq_k)}{S}
\end{equation}
wobei:
\begin{itemize}[itemsep=2pt]
\item $\omega \in \{1, \ldots, 256\}$ der ganzzahlige Relaxationsparameter ist,
\item $S = 256$, sodass die Division als Rechts-Shift $\gg 8$ implementiert wird.
\end{itemize}
\end{definition}

\begin{remark}[Physikalische Interpretation von $\omega$]~
\begin{itemize}[itemsep=4pt]
\item $\omega = 256$ ($= S$): \textbf{Sofortige Relaxation.}
Agent springt in einem Schritt ins Gleichgewicht.
Resultat: abrupte Richtungswechsel, ruckartige Bewegung.
Analog zu einer d\"unnfl\"ussigen Substanz (Wasser).

\item $\omega = 128$ ($= S/2$): \textbf{Halbe Relaxation pro Schritt.}
Agent n\"ahert sich dem Gleichgewicht in $\sim 2$ Schritten.
Resultat: glatte, flie\ss ende Bewegung. \textbf{Empfohlener Standardwert.}

\item $\omega = 32$ ($= S/8$): \textbf{Langsame Relaxation.}
Agent braucht $\sim 8$ Schritte f\"ur eine Richtungs\"anderung.
Resultat: sehr tr\"age, vorausschaubare Bewegung.
Analog zu einer dickfl\"ussigen Substanz (Honig).

\item $\omega \to 0$: Keine Relaxation. Agent ignoriert das Gleichgewicht.
Bewegung friert ein. \textbf{Nicht empfohlen} (f\"uhrt zu Deadlocks).
\end{itemize}
\end{remark}

\begin{proposition}[Stabilit\"atsbedingung]
F\"ur numerische Stabilit\"at muss gelten:
\begin{equation}
0 < \omega \leq 2S = 512
\end{equation}
Der stabile und praktisch sinnvolle Bereich ist $\mathbf{32 \leq \omega \leq 256}$.
\end{proposition}

\noindent\textbf{Implementierung:}

\begin{algorithm}[H]
\SetAlgoLined
\KwIn{$f[0..16]$, $\feq[0..16]$, $\omega$}
\KwOut{$f^*[0..16]$}
\For{$k \leftarrow 0$ \KwTo $16$}{
    $\Delta \leftarrow f[k] - \feq[k]$\;
    $f^*[k] \leftarrow f[k] - (\omega \cdot \Delta) \gg 8$\;
    \tcp{Nicht-Negativit\"at sicherstellen}
    \lIf{$f^*[k] < 0$}{$f^*[k] \leftarrow 0$}
}
\caption{Integer-Relaxation (BGK-Schema)}
\end{algorithm}

\textbf{Kosten:} $17 \times (1\text{ Sub} + 1\text{ Mul} + 1\text{ Shift} + 1\text{ Vgl}) = 68$ Integer-Zyklen pro Agent.

% ============================================================
\section{Positionsupdate (Streaming)}\label{sec:streaming}
% ============================================================

In der Fluid-LBM bewegen sich Verteilungswerte zu Nachbargitterpunkten.
F\"ur Agenten im kontinuierlichen Raum gibt es kein Gitter.
Stattdessen bestimmt der Verteilungsvektor nach der Relaxation
die \textbf{Bewegungsrichtung und -geschwindigkeit} des Agenten:

\begin{definition}[Positionsupdate]
Der Bewegungsimpuls nach Relaxation ist:
\begin{equation}
\mathbf{j}^* = \sum_{k=1}^{16} f^*_k \cdot \mathbf{c}_k
\end{equation}
Die neue Position ergibt sich als:
\begin{equation}
\mathbf{x}_A(t+\Delta t) = \mathbf{x}_A(t) + \frac{\mathbf{j}^* \cdot \Delta t}{\rho \cdot S}
\end{equation}
\end{definition}

In Integer-Arithmetik:
\begin{align}
j^*_x &= \sum_{k=1}^{16} f^*_k \cdot c_{k,x} \\
j^*_y &= \sum_{k=1}^{16} f^*_k \cdot c_{k,y} \\
x_{\mathrm{neu}} &= x_{\mathrm{alt}} + (j^*_x \cdot \Delta t) \gg 16 \label{eq:xneu}\\
y_{\mathrm{neu}} &= y_{\mathrm{alt}} + (j^*_y \cdot \Delta t) \gg 16 \label{eq:yneu}
\end{align}

Der Shift $\gg 16$ in (\ref{eq:xneu})--(\ref{eq:yneu}) kombiniert die Division durch $\rho \approx 256$ ($\gg 8$)
und durch $S = 256$ ($\gg 8$) in einer einzigen Operation.

\textbf{Kosten:} $16 \times 2$ Multiplikationen + 2 Shifts + 2 Additionen = 36 Zyklen.

% ============================================================
\section{Hindernisbehandlung}\label{sec:bounce}
% ============================================================

In der LBM werden Wandrandbedingungen durch \textit{Bounce-Back} realisiert:
Verteilungswerte, die auf ein Hindernis treffen, werden in die antiparallele Richtung reflektiert.
Wir \"ubertragen dieses Prinzip direkt.

\begin{definition}[Compass16-Bounce-Back]
Wenn Agent $A$ ein Hindernis in Richtung $k$ detektiert:
\begin{equation}
f_k \leftrightarrow f_{\bar{k}}
\end{equation}
wobei $\bar{k} = ((k-1+8) \bmod 16) + 1$ der antiparallele Index ist (Abschnitt~2).
\end{definition}

\textbf{Hindernis-Detektion via DSM:} Ein Hindernis in Richtung $k$ liegt vor, wenn
ein Objekt $j$ existiert mit:
\begin{enumerate}[itemsep=2pt]
\item $\compass(\mathbf{x}_j - \mathbf{x}_A) = k$ \quad (Objekt liegt in Richtung $k$)
\item $d^2(\mathbf{x}_j, \mathbf{x}_A) < d^2_{\mathrm{sicher}}$ \quad (Objekt ist zu nah)
\end{enumerate}
Beide Bedingungen sind reine Integer-Operationen --- identisch zum DSM-Matching.

\textbf{Unterscheidung statisch/dynamisch:}
Statische Hindernisse (Regalw\"ande, Mauern) haben eine feste Position.
Dynamische Hindernisse (andere Agenten) bewegen sich.
Die Lighthouse-Signaturen aus DSM erm\"oglichen diese Unterscheidung:
Statische Objekte haben konstante Lighthouse-Signaturen \"uber die Zeit,
dynamische Objekte haben ver\"anderliche. Die Abfrage kostet 2 Integer-Vergleiche pro Leuchtturm.

% ============================================================
\section{Vollst\"andiger Algorithmus}
% ============================================================

\begin{algorithm}[H]
\SetAlgoLined
\KwIn{$N$ Agenten mit Positionen $\mathbf{x}^{(A)}$, Verteilungen $\mathbf{f}^{(A)}$, Zielen $\mathbf{x}^{(A)}_{\mathrm{Ziel}}$}
\KwOut{Neue Positionen und Verteilungen}
\BlankLine
\ForEach{Agent $A$}{
    \tcp{\textbf{Phase 1: Lokale Wahrnehmung} (DSM-Kern)}
    \ForEach{Nachbar $j \in \mathcal{N}(A)$}{
        $d^2_j \leftarrow (\Delta x)^2 + (\Delta y)^2$ \tcp*{5 Zyklen}
        $\kappa_j \leftarrow \compass(\Delta x, \Delta y)$ \tcp*{12 Zyklen}
    }
    \BlankLine
    \tcp{\textbf{Phase 2: Sollbewegung berechnen} (Gl.~\ref{eq:utarget})}
    $\kappa_Z \leftarrow \compass(\mathbf{x}_{\mathrm{Ziel}} - \mathbf{x}_A)$\;
    $u_x \leftarrow \alpha \cdot c_{\kappa_Z,x}$ \tcp*{Navigation}
    $u_x \leftarrow u_x - \beta \cdot {\textstyle\sum_{j:\, d^2_j < d^2_{\min}}} (d^2_{\min}/d^2_j)\, c_{\kappa_j,x}$ \tcp*{Separation}
    $u_x \leftarrow u_x + \gamma \cdot {\textstyle\sum_{j:\, d^2_j > d^2_{\max}}} (d^2_j/d^2_{\max})\, c_{\kappa_j,x}$ \tcp*{Koh\"asion}
    $u_y \leftarrow$ analog f\"ur $y$-Komponente\;
    \BlankLine
    \tcp{\textbf{Phase 3: Gleichgewicht berechnen} (Gl.~\ref{eq:feq})}
    \For{$k \leftarrow 0$ \KwTo $16$}{
        $\mathrm{dot}_k \leftarrow (c_{k,x} \cdot u_x + c_{k,y} \cdot u_y) \gg 8$\;
        $\feq_k \leftarrow (\hat{w}_k \cdot (\rho + (\kappa \cdot \mathrm{dot}_k) \gg n_\kappa)) \;/\; W$\;
    }
    \BlankLine
    \tcp{\textbf{Phase 4: Relaxation} (Gl.~\ref{eq:collision})}
    \For{$k \leftarrow 0$ \KwTo $16$}{
        $f^*_k \leftarrow f_k - (\omega \cdot (f_k - \feq_k)) \gg 8$\;
        \lIf{$f^*_k < 0$}{$f^*_k \leftarrow 0$}
    }
    \BlankLine
    \tcp{\textbf{Phase 5: Hindernisbehandlung} (Bounce-Back)}
    \For{$k \leftarrow 1$ \KwTo $16$}{
        \If{Hindernis in Richtung $k$ mit $d^2 < d^2_{\mathrm{sicher}}$}{
            Tausche $f^*_k \leftrightarrow f^*_{\bar{k}}$\;
        }
    }
    \BlankLine
    \tcp{\textbf{Phase 6: Positionsupdate}}
    $j^*_x \leftarrow \sum_k f^*_k \cdot c_{k,x}$; \quad
    $j^*_y \leftarrow \sum_k f^*_k \cdot c_{k,y}$\;
    $x_{\mathrm{neu}} \leftarrow x_{\mathrm{alt}} + (j^*_x \cdot \Delta t) \gg 16$\;
    $y_{\mathrm{neu}} \leftarrow y_{\mathrm{alt}} + (j^*_y \cdot \Delta t) \gg 16$\;
}
\caption{Compass16-LBM: Vollst\"andiger Zeitschritt}
\end{algorithm}

% ============================================================
\section{Kostenanalyse}
% ============================================================

\begin{center}
\begin{tabular}{@{}lrrr@{}}
\toprule
\textbf{Phase} & \textbf{Operationen} & \textbf{Zyklen/Agent} & \textbf{FP-Ops} \\
\midrule
1. Wahrnehmung ($k{=}3$ Nachbarn) & $3 \times 17$ & $51$ & 0 \\
2. Sollbewegung & $3 + 2k$ Mul/Add & $\sim 20$ & 0 \\
3. Gleichgewicht (17 Werte) & $17 \times 5$ & $85$ & 0 \\
4. Relaxation (17 Werte) & $17 \times 4$ & $68$ & 0 \\
5. Hindernisbehandlung & $\leq 16$ Swaps & $\leq 16$ & 0 \\
6. Positionsupdate & $16 \times 2 + 4$ & $36$ & 0 \\
\midrule
\textbf{Gesamt pro Agent pro Zeitschritt} & & $\mathbf{\sim 276}$ & $\mathbf{0}$ \\
\bottomrule
\end{tabular}
\end{center}

\noindent Hochrechnungen f\"ur ATmega328 (16\,MHz, 1 Zyklus = 62{,}5\,ns):

\begin{center}
\begin{tabular}{@{}rrrl@{}}
\toprule
$N$ (Agenten) & Zyklen/Schritt & Zeit/Schritt & Update-Rate \\
\midrule
10 & 2\,760 & 0{,}17\,ms & $>$1000\,Hz \\
100 & 27\,600 & 1{,}7\,ms & $\sim$580\,Hz \\
500 & 138\,000 & 8{,}6\,ms & $\sim$116\,Hz \\
1\,000 & 276\,000 & 17{,}3\,ms & $\sim$58\,Hz \\
\bottomrule
\end{tabular}
\end{center}

Zum Vergleich: Ein Reynolds-Flocking-Algorithmus mit $\sqrt{\cdot}$ und $\arctan$
ben\"otigt $\sim$2\,000--3\,000 Zyklen/Agent auf derselben Hardware ($\sim$8--10$\times$ langsamer).

% ============================================================
\section{Parameterempfehlungen}
% ============================================================

\begin{center}
\begin{tabular}{@{}lllp{5.5cm}@{}}
\toprule
\textbf{Parameter} & \textbf{Symbol} & \textbf{Empfehlung} & \textbf{Bedeutung} \\
\midrule
Skalierung & $S$ & 256 & Festkomma-Aufl\"osung (Shift-kompatibel) \\
Reaktivit\"at & $\kappa$ & 4 & Wie stark beeinflusst die Sollbewegung das Gleichgewicht \\
Relaxation & $\omega$ & 128 & Gl\"atte der Richtungs\"anderung \\
Navigation & $\alpha$ & 128 & St\"arke der Zielnavigation \\
Separation & $\beta$ & 256 & St\"arke der Kollisionsvermeidung \\
Koh\"asion & $\gamma$ & 64 & St\"arke des Formationserhalts \\
Sicherheitsabstand$^2$ & $d^2_{\min}$ & 400 & Minimaler Abstand ($\approx$20 Einheiten) \\
Koh\"asionsabstand$^2$ & $d^2_{\max}$ & 10\,000 & Maximaler Abstand ($\approx$100 Einheiten) \\
Leuchtturm-Anzahl & $m$ & 3 & Globale Referenzpunkte \\
\bottomrule
\end{tabular}
\end{center}

Alle Parameter sind Zweierpotenzen oder lassen sich als solche w\"ahlen,
sodass Multiplikationen konsequent durch Shift-Operationen ersetzt werden k\"onnen.

% ============================================================
\section{Emergente Eigenschaften}
% ============================================================

Durch die LBM-Struktur ergeben sich folgende Verhaltensweisen \textit{emergent}
--- sie werden nicht explizit programmiert, sondern entstehen aus den lokalen Regeln:

\textbf{1. Flie\ss verhalten um Hindernisse.}
Bei $\omega \approx 128$ und $\beta > \alpha$ weichen Agenten
Hindernissen aus wie eine Fl\"ussigkeit, die um einen Stein flie\ss t.
Die Verteilungswerte verteilen sich automatisch auf die freien Compass-Sektoren.
Es muss kein Ausweichpfad berechnet werden --- er ergibt sich aus der Relaxation.

\textbf{2. Automatischer Abstandsausgleich.}
Hohe lokale Agentendichte erzeugt starke Separationsbeitr\"age,
die Agenten auseinandertreiben.
Niedrige Dichte erzeugt Koh\"asionsbeitr\"age, die Agenten zusammenziehen.
Der Gleichgewichtsabstand stellt sich selbstorganisiert ein.

\textbf{3. Geordneter Gegenverkehr in Engstellen.}
Bei schmalen Passagen (z.B.\ Lagergassen) erzeugt der Bounce-Back
an den W\"anden eine asymmetrische Verteilung,
die nat\"urlichen Gegenverkehr mit impliziter Spurbildung hervorbringt.

\textbf{4. Deadlock-Freiheit.}
Die Ruhe-Verteilung $f_0$ wird durch die Relaxation
st\"andig zum Gleichgewicht gezogen.
Da das Gleichgewicht stets eine Zielrichtung enth\"alt ($\mathbf{u}_{\mathrm{Ziel}} \neq \mathbf{0}$,
solange das Ziel nicht erreicht ist),
kann kein Agent dauerhaft stehen bleiben. Das System bewegt sich immer.

\textbf{5. Gummiseil-Verhalten.}
Die Kombination aus Separation ($d^2 < d^2_{\min}$: Absto\ss ung)
und Koh\"asion ($d^2 > d^2_{\max}$: Anziehung)
mit stetigem \"Ubergang erzeugt das ``virtuelle Gummiseil'':
Agenten werden bei \"Uberdehnung angezogen und bei Kompression abgesto\ss en,
mit einer ``Ruhelänge'' im Bereich $[d^2_{\min}, d^2_{\max}]$.

% ============================================================
\section{Offene Fragen}
% ============================================================

\begin{enumerate}[itemsep=4pt]
\item \textbf{Isotropie der Gewichte.}
Die einfache Wahl $w_r = 1/24$ f\"ur alle 16 Richtungen
ist eine N\"aherung erster Ordnung.
Exakte Isotropie 5.\ Ordnung erfordert m\"oglicherweise
unterschiedliche Gewichte f\"ur die drei Geschwindigkeitsklassen
(Kardinal: $|\mathbf{c}| = 256$, Diagonal: $|\mathbf{c}| = 362$, Zwischenrichtungen: $|\mathbf{c}| = 256$).
Numerische Validierung durch Simulation steht aus.

\item \textbf{Integer-Division in der Separation.}
Die Division $d^2_{\min}/d^2_j$ ist nicht direkt als Shift darstellbar.
Drei L\"osungsans\"atze:
(a)~Lookup-Table f\"ur $256/d^2$ (256 Eintr\"age, 512 Bytes ROM),
(b)~st\"uckweise lineare Approximation in 4 Segmenten,
(c)~Ersetzen durch $(d^2_{\min} - d^2_j)$ (lineare statt inverse Absto\ss ung).
Option (c) eliminiert die Division vollst\"andig, \"andert aber das Absto\ss ungsprofil.

\item \textbf{3D-Erweiterung.}
F\"ur fliegende Agenten (Drohnen) wird eine sph\"arische Kompassrose ben\"otigt.
M\"ogliche Ans\"atze: Compass16 $\times$ 5 Elevationsklassen = 80+1 Richtungen (162 Bytes/Agent),
oder Compass32 $\times$ 3 Elevationsklassen = 96+1 Richtungen (194 Bytes/Agent).
Machbar auf 32-Bit-MCUs, eng auf 8-Bit-Systemen.

\item \textbf{Stabilit\"atsanalyse.}
Unter welchen Bedingungen konvergiert das System zu einer stabilen Formation?
Die LBM-Theorie bietet hierf\"ur analytische Werkzeuge
(Chapman-Enskog-Entwicklung), die auf das Schwarmmodell \"ubertragbar sein k\"onnten.
Eine rigorose Analyse steht aus.

\item \textbf{Validierung gegen Baselines.}
Vergleich mit: (a)~Standard-Reynolds-Flocking (mit FPU),
(b)~ORCA (Optimal Reciprocal Collision Avoidance),
(c)~zentralisiertem MAPF (CBS-Algorithmus).
Metriken: Kollisionsrate, mittlere Ankunftszeit, Zyklen pro Zeitschritt.
\end{enumerate}

% ============================================================
\section*{Referenzen}
% ============================================================

\begin{itemize}[itemsep=2pt,leftmargin=*]
\item Balban, D. (2026). \textit{Trigonometry-Free TSP Heuristics for Embedded Systems:
A Compass-Based Matching Approach.} Working Paper. DOI: 10.5281/zenodo.19096781.
\item Bhatnagar, P.L., Gross, E.P., Krook, M. (1954). \textit{A model for collision processes
in gases.} Physical Review, 94(3), 511--525.
\item Reynolds, C.W. (1987). \textit{Flocks, herds and schools: A distributed behavioral model.}
SIGGRAPH Computer Graphics, 21(4), 25--34.
\item Qian, Y.H., d'Humi\`eres, D., Lallemand, P. (1992). \textit{Lattice BGK models
for Navier-Stokes equation.} Europhysics Letters, 17(6), 479--484.
\item van den Berg, J., et al. (2011). \textit{Reciprocal n-Body Collision Avoidance.}
Robotics Research, STAR 70, 3--19.
\item Sharon, G., et al. (2015). \textit{Conflict-Based Search for Optimal Multi-Agent Pathfinding.}
Artificial Intelligence, 219, 40--66.
\end{itemize}

\end{document}
